\documentclass[10pt]{article}
\usepackage[letterpaper,margin=.38in]{geometry}
\usepackage{amsmath,amssymb,multicol}
\pdfinfoomitdate=1\pdftrailerid{}
\pagestyle{empty}\setlength{\parindent}{0pt}\setlength{\parskip}{2pt}
\setlength{\columnsep}{13pt}\setlength{\jot}{0pt}
\setlength{\abovedisplayskip}{3pt}\setlength{\belowdisplayskip}{3pt}
\let\bigsum\sum\renewcommand{\sum}{\bigsum\nolimits}
\newcommand{\fl}[1]{\left\lfloor#1\right\rfloor}
\newcommand{\rg}[2]{\{#1,\ldots,#2\}}
\begin{document}
\begin{center}
{\large\bfseries A thirteen-variable consistency equation}\\[-1pt]
{\small Positive integer unknowns; exact degree 56; an 86-operation evaluator.}
\end{center}\vspace{-8pt}
\begin{multicols}{2}
\fontsize{8.6}{9.1}\selectfont
% BEGIN VERIFIED DISPLAY
The only unknowns are the thirteen distinct positive integers
\[
J,F,\alpha,f,g,T,X,Y,\xi,\eta,\zeta,\upsilon,\gamma\in\mathbb Z_{>0}.
\]
All other symbols below are polynomial abbreviations:
\[
\begin{gathered}
b=\beta J,\quad q=1+b,\quad C=q-F-\alpha,\quad U=Yq^2,\\
E=XY,\quad k=\eta+\zeta,\quad \ell=EYk,\quad c=\eta+Yk,\\
a=E+Y,\quad H=3+4a,\quad\Delta=H+a^2,\quad L=X+ac,\\
M=q^3b,\quad W=C(X+Kq)-Fq,\\
R=[\alpha+(q-F)b](q^2-1)+J(\mu+\omega q),\\
Z=M(\ell+k+\xi-kW)+Yb-Wq^3,\\
D=UL+H(U\gamma+Eq+R+1-k),\\
I=c^2g,\quad Q=\Delta(f^2-1),\quad V=f^2(Tc-R)-c.
\end{gathered}
\]
The single Diophantine equation is
\[
\boxed{\begin{aligned}
&[1+Z^2-M^2(1+\ell(\ell+k))]\\
&\quad\cdot[1+D^2-U^2(1+\Delta c^2)]\\
&\quad\cdot[\upsilon^2-Q(\upsilon^2-V^2)]
[1+I^2-Q]=1.
\end{aligned}}
\]
% END VERIFIED DISPLAY
\textbf{The four fixed integers.}
Every index, set, sum, product and floor below defines a constant, not
another unknown. Put
\[
\begin{gathered}
\rho=2^{47},\quad B=\rho^{2011},\quad\beta=B-1,\\
G_m(x)=\sum_{j=0}^{m-1}x^j,\quad
\Lambda=\rho^{1005}.
\end{gathered}
\]
\begin{align*}
\mu={}&\beta-G_{951}(\rho)\\
&-\rho^{960}G_{26}(\rho)-\rho^{1004},\\
\omega={}&\beta+8G_{917}(\rho)-4+(2^{45}-4)\Lambda\\
&+\rho^{960}G_9(\rho)+\rho^{987}G_8(\rho)
  +\rho^{1004},\\
K={}&\rho^{18}+\rho^{1004}+\rho^{2020}\\
&+2\Lambda\Phi(1/\rho),
\end{align*}
where
\[
\begin{gathered}
\chi_i=1+2^{10}(1+w_i+2^8n_i\\
\qquad+2^{17}p_i+2^{25}s_i)\quad(0\le i\le941),\\
\pi(0)=0,\quad\pi(i)=1+\sum_{j=1}^{941}
\fl{\frac{2^{44}-1+\chi_j-\chi_i}{2^{44}}}\quad(1\le i\le941).
\end{gathered}
\]
\[
\begin{split}
\Phi(x)={}&\sum_{i=0}^{941}\chi_i x^{\pi(i)}\\
&+2^{10}x^{969}(x^9+2^{17})G_8(2x)\\
&+2^{18}x^{942}(x^{18}+2^{17})G_9(2x).
\end{split}
\]
The rational intermediate in K cancels exactly; all four results are positive integers.

\textbf{The 942 indexed fixed tiles.}
Use the actions $u_j,v_j,d_j$ decoded below. For $d=0,1$, let
\[
\begin{gathered}
h(q,b)=2q+b+2,\quad
\mathcal S_d=\{u_j:d_j=d,\ 0\le j<430\},\\
r_{d,q}=205-199d+|\mathcal S_d\cap\rg0{q-1}|.
\end{gathered}
\]
An interval ending at $-1$ is empty. The table defines $t_i$ for
$0\le i\le941$; all index ranges are fixed:
\[
\setlength{\arraycolsep}{1.4pt}\begin{array}{c|l}
\text{index }i&\text{tile }t_i=(w_i,n_i,p_i,s_i)\\\hline
0&(2,435,3,1)\\
j+1\ (j=2q+b<430)&((1-d_j)r_{0,u_j},h(q,b),d_jr_{1,u_j},v_j)\\
433-2a+b\ (a,b=0,1)&(0,432a+b,0,432a+b)\\
\substack{423+2r_{d,q}+b\\q\in\mathcal S_d,\ d,b=0,1}&(dr_{d,q},b,(1-d)r_{d,q},h(q,b))\\
923+2a+b\ (a,b=0,1)&(0,432a+b,0,438)\\
927&(2,434,2,0)\\
928&(3,436,4,4)\\
929&(4,436,4,0)\\
496+t\ (t=434,435,436)&(1,438,1,t)\\
933+3a+b\ (a,b=0,1,2)&(a^2,437,b,437)
\end{array}
\]
\textbf{A closed decoder for the 430 actions.}
Put $D_b(A,k)=\fl{A/b^k}\bmod b$, where
$a\bmod b=a-b\fl{a/b}$. For $0\le q<215$, $b=0,1$, $j=2q+b$:
\[
\begin{gathered}
\mathcal N=\mathcal A\bmod2^{636},\quad
\mathcal T=\fl{\mathcal A/2^{636}},\\
\nu_j=D_2(\mathcal N,j),\quad \kappa_j=(\sigma_j+2)^{1-\nu_j}\quad(0\le j<430),\\
\sigma_j=\sum_{k=0}^{j-1}\nu_k,\quad
\Pi_j=\prod_{k=0}^{j-1}\kappa_k\quad(0\le j\le430),\\
m_q=D_2(\mathcal N,430+q),\quad \tau_q=\sum_{k=0}^{q-1}m_k,\\
F_q=14+m_q(D_{16}(\fl{\mathcal T/\Pi_{430}},\tau_q)-14),\\
u_j=(\sigma_{j+1}-(\fl{\mathcal T/\Pi_j}\bmod\kappa_j))\bmod216,\\
v_j=D_2(F_q,2b),\quad d_j=D_2(F_q,2b+1).
\end{gathered}
\]
Empty sums/products are $0/1$. Here $(v_0,d_0,u_0)=(1,1,1)$; the constructor tile
uses this actual first target, including its state renaming.

\textbf{The integer $\mathcal A$ in base 10.}
Its complete decimal expansion is the concatenation of these lines:\par
\begingroup\fontsize{7.7}{8.2}\selectfont\ttfamily
12889260410359561746736540743515845134244260640012702033212621\par
95398074180784058571518488282119861023093355389406939414042892\par
97455773992193784795705001618777025966369137954859764532832722\par
75411843410809038713502633204617333246733616940964578272360901\par
56186216508504953649192752154519268024509952082950238865468307\par
22999632482769253980274091217460707829023970817244790340704449\par
09085581271192071656170732004937337082171141082486384269072884\par
72778087419663417942892520175630649792329265834490253035577533\par
96449689808375788457488971017043689875433162972390177936123032\par
89025657757489526342497438772045714099661532481123050235898973\par
27443609903800541199351405985754278307227777584454508422717998\par
334113564409335346074892680328437\par
\endgroup
\normalfont\fontsize{7.7}{8.2}\selectfont
The fixed integer $\mathcal A$ above is written in base 10. A Lean theorem
establishes the existence of an HA derivation of both directions between
this no-solution sentence and the formalized consistency of ZFC. The theorem
was independently replayed on a checked pre-merge source snapshot; historical
HA, ZF, Foundation and package inputs remain inherited boundaries. A later
source refresh is in progress. Exact degree 56 is certified by integer-polynomial expansion. Its 87-operation
parent has degree 51: one fewer gate trades against a higher degree. Neither
count is a minimality claim.
\end{multicols}
\clearpage
\begin{center}{\large\bfseries Exact 86-operation circuit}\end{center}
Each numbered line defines one gate over the integers; the inputs are the
thirteen unknowns and the four fixed integers above. The circuit uses the
ASCII names \texttt{alpha, xi, eta, zeta, upsilon, gamma, beta, mu, omega,}
\texttt{ell, Delta} for their displayed counterparts; \texttt{g} is the
unknown named $g$. Its last gate is zero
exactly when the displayed equation holds. Intermediate gate names below are
local names, not additional unknowns.
\begin{multicols}{3}
\raggedright\ttfamily\fontsize{7.2}{8.4}\selectfont
1. b = J * beta\\
2. q = 1 + b\\
3. q2 = q * q\\
4. q3 = q * q2\\
5. U = Y * q2\\
6. qF = q - F\\
7. C = qF - alpha\\
8. E = X * Y\\
9. k = eta + zeta\\
10. kY = Y * k\\
11. c = eta + kY\\
12. c2 = c * c\\
13. a = E + Y\\
14. g15 = 4 * a\\
15. H = 3 + g15\\
16. g17 = a * a\\
17. Delta = H + g17\\
18. g19 = a * c\\
19. L = X + g19\\
20. ell = E * kY\\
21. M = q3 * b\\
22. g25 = K * q\\
23. g26 = X + g25\\
24. g27 = C * g26\\
25. g28 = F * q\\
26. W = g27 - g28\\
27. g30 = ell + k\\
28. g31 = g30 + xi\\
29. root\_skew = k * W\\
30. root\_inner = g31 - root\_skew\\
31. root\_scaled = M * root\_inner\\
32. g33 = Y * b\\
33. root\_btq3 = W * q3\\
34. root\_sum = root\_scaled + g33\\
35. Z = root\_sum - root\_btq3\\
36. g38 = Z * Z\\
37. g39 = M * M\\
38. g40 = ell * g30\\
39. g41 = 1 + g40\\
40. g42 = g39 * g41\\
41. A0 = g38 - g42\\
42. g44 = qF * b\\
43. rankseed = alpha + g44\\
44. g46 = q2 - 1\\
45. g47 = rankseed * g46\\
46. g48 = omega * q\\
47. g49 = mu + g48\\
48. g50 = J * g49\\
49. R = g47 + g50\\
50. g52 = k - R\\
51. r0 = g52 - 1\\
52. g54 = U * gamma\\
53. g55 = E * q\\
54. g56 = g54 + g55\\
55. shift = g56 - r0\\
56. g58 = U * L\\
57. g59 = H * shift\\
58. D = g58 + g59\\
59. g61 = D * D\\
60. g62 = U * U\\
61. g63 = Delta * c2\\
62. g64 = 1 + g63\\
63. g65 = g62 * g64\\
64. A1 = g61 - g65\\
65. f2 = f * f\\
66. rho = f2 - 1\\
67. Q = Delta * rho\\
68. I = c2 * g\\
69. g71 = I * I\\
70. g72 = g71 - Q\\
71. strong = 1 + g72\\
72. g74 = T * c\\
73. g75 = g74 - R\\
74. g76 = f2 * g75\\
75. V = g76 - c\\
76. y2 = upsilon * upsilon\\
77. g80 = V * V\\
78. g81 = y2 - g80\\
79. g82 = Q * g81\\
80. aux = y2 - g82\\
81. first = 1 + A0\\
82. main = 1 + A1\\
83. g86 = first * main\\
84. g87 = aux * g86\\
85. g88 = strong * g87\\
86. P = g88 - 1\\
\end{multicols}
\end{document}
